%Tex root = ProbabilitàEStatistica.tex

\chapter{Correlazione tra variabili e regressione lineare}

\section{Correlazione tra variabili}
Talvolta, per ogni individuo vengono misurati più caratteri: peso, altezza, sessp, reddito, ecc.\\[0.2ex]
Consideriamo, ora, il caso dei due cartteri quantitativi e supponiamo che siano sotto forma di coppie:
\begin{equation*}
(x_1,y_1), (x_2,y_2), \dots, (x_n,y_n)
\end{equation*}

Ogni coppia fa riferimento a un individuo.

\medskip
In un primo approcio grafico, si possono disegnare sul piano tutti i punti di cordinate $(x_i,y_i)$ e vedere se essi tendono a disporsi secondo un andamento regolare.\\[0.2ex]
QUesto tipo di rappresentazione grafica si chiama \textcolor{viola}{\textbf{scatterplot}} o \textcolor{viola}{\textbf{digramma di dispersione}}.

%foto 4 scatterplot

L'idea è quella di trovare la curva (retta, parabola ecc.), se esiste, che meglio descriva l'andamento dei dati per poi utilizzarla per "stimare" il valore di un carattere conoscendo l'altro.

\Definizione{covarianza}{
Supponiamo di avere N osservazioni congiunte di due variabili X e Y:
\begin{equation*}
    {(x_1,y_1), \dots, (x_n,y_n)}
\end{equation*}

Si chiama \textbf{covarianza} tra X e Y la quantità:
\begin{equation*}
    s_{xy} = \frac{1}{N}\sum_{i=1}^{N}(x_i - \bar{X})(y_i - \bar{Y}) = \bar{XY} - \bar{X} \cdots \bar{Y}
\end{equation*}
}

Si dice che le variabili X e Y sono:
\begin{itemize}[label=$\triangleright$]
\item \textbf{direttamente correlate} se $s_{xy} > 0$
\smallskip
\item \textbf{inversamente correlate} se $s_{xy} < 0$
\smallskip
\item \textbf{incorrelate} se $s_{xy} = 0$
\end{itemize}

\Definizione{coefficiente di correlazione}{
Si chiama \textbf{coefficiente di correlazione} di X e Y la quantità:
\begin{equation*}
    \rho_{xy} = \frac{s_{xy}}{s_xs_y}
\end{equation*}

ovvero, la covarianza divisa per il prodotto delle deviazioni standard.

\noindent\rule{\linewidth}{0.4pt}
\textbf{N.B.} coefficiente di correlazione è definito solo se $s_x \neq 0$ e $s_y \neq 0$ (se entrambi non sono costanti).
}

\vario{Proprietà del coefficiente di correlazione campionaria}{
Di seguito diamo alcune delle propreità del coefficiente di correlazione campionaria $\rho_{xy}$:
\begin{itemize}[label=$\circ$]
    \item sempre $-1 \leq \rho_{xy} \leq 1$
    \smallskip
    \item se per opportune costanti a e b (b > 0), sussite la relazione lineare:
    \begin{equation*}
    y_i = a + bx, \quad \forall i = 1, 2, \dots n
    \end{equation*}

    allora $\rho_{xy} = 1$.
    \smallskip
    \item se per opportune costanti a e b (b < 0), sussite la relazione lineare:
    \begin{equation*}
    y_i = a + bx, \quad \forall i = 1, 2, \dots n
    \end{equation*}

    allora $\rho_{xy} = -1$.
    \smallskip
    \item se $\rho_{xy}$ è il coefficiente di correlazione del campione $(x_i,y_i), \quad i = 1, \dots, n$, allora lo è anche per il campione:
    \begin{equation*}
    (a + bx_i, c + dy_i), \quad \forall i = 1, 2, \dots n
    \end{equation*}

    purchè le costanti d e d abbiano lo stesso segno.
\end{itemize}
}

\section{Regressione lineare}
$\dots$ verrà trattato più avanti